\chapter{Ética na Pesquisa e Responsabilidade Acadêmica}

A ética na pesquisa envolve respeito às pessoas, aos animais, às fontes, aos dados, às instituições e à autoria. Ela deve orientar o desenho do estudo desde o projeto, e não ser acrescentada apenas como declaração formal na versão final.

\section{Pesquisas com seres humanos}

Entrevistas, questionários, grupos focais, observação, intervenções e acesso a prontuários, imagens, gravações ou dados pessoais podem caracterizar pesquisa envolvendo seres humanos. O enquadramento deve ser verificado com o orientador e com a instância institucional competente antes do recrutamento ou da coleta.

Quando a apreciação pelo Sistema CEP/Conep for aplicável, o protocolo deve ser apresentado com a documentação exigida e a pesquisa somente pode começar após a aprovação ética. Nos TCCs de graduação abrangidos pela Resolução CNS nº 510/2016, o registro no Comitê de Ética em Pesquisa ocorre sob responsabilidade do orientador. Pesquisa vinculada a projeto do orientador já aprovado pode exigir emenda, conforme o caso.

Nem toda atividade acadêmica com pessoas possui o mesmo enquadramento. Atividades exclusivamente educacionais, levantamentos sem finalidade de pesquisa e uso de informações de acesso público, por exemplo, devem ser analisados de acordo com a regulamentação aplicável; o estudante não deve presumir dispensa por conta própria.

\caixaatencao{Consulte o orientador e a instância competente antes de recrutar participantes, acessar prontuários ou iniciar qualquer coleta. O manual não autoriza dispensa de apreciação ética.}

\subsection{Documentos e autorizações}

Conforme o desenho da pesquisa, podem ser necessários projeto detalhado, instrumentos de coleta, termos de consentimento ou assentimento, autorização do local ou da instituição, autorização para acesso a registros e parecer ético. Esses documentos têm funções diferentes: a autorização institucional não substitui a aprovação ética, e nenhuma delas substitui o processo de consentimento dos participantes quando ele for aplicável.

\section{Pesquisas com animais e material biológico}

Atividades de ensino ou pesquisa científica com animais devem observar a Lei nº 11.794/2008, as normas do Conselho Nacional de Controle de Experimentação Animal e os procedimentos da Comissão de Ética no Uso de Animais. Quando exigida, a autorização da CEUA deve ser obtida antes do início da atividade.

Pesquisas com material biológico humano, dados clínicos ou procedimentos de saúde podem exigir cuidados adicionais. O estudante deve verificar o fluxo institucional e a regulamentação específica da área ainda na fase de projeto.

\section{Consentimento, assentimento e privacidade}

O consentimento livre e esclarecido é um processo de comunicação. O participante deve compreender objetivos, procedimentos, riscos, benefícios, voluntariedade, formas de registro e uso das informações. Quando houver crianças, adolescentes ou pessoas sem plena capacidade de consentir, deve-se verificar também a necessidade de assentimento e de autorização do responsável.

A forma de registro do consentimento depende do contexto e das normas aplicáveis. Assinatura em papel não é a única possibilidade em todos os estudos, mas qualquer alternativa deve ser justificada e aprovada quando a pesquisa estiver sujeita à apreciação ética.

\section{Proteção, anonimização e gestão dos dados}

O projeto deve definir quais dados são realmente necessários, quem terá acesso, onde serão armazenados, por quanto tempo serão mantidos e como serão eliminados ou compartilhados. Dados pessoais e sensíveis exigem medidas proporcionais ao risco, conforme a legislação e as orientações institucionais.

O plano de gestão deve registrar finalidade, base legal a ser confirmada institucionalmente, responsáveis, controle de acesso, cópias de segurança, retenção, descarte e condições de compartilhamento. Dados não devem ser enviados a serviços externos ou ferramentas de inteligência artificial sem autorização e avaliação do risco de exposição.

Anonimização e pseudonimização não são sinônimos. Na pseudonimização, os registros ainda podem ser associados a uma pessoa mediante informação adicional mantida separadamente. A anonimização busca impedir a reidentificação por meios razoáveis, mas deve ser tratada como processo baseado em risco, e não como garantia absoluta.

Imagens, voz, gravações e combinações de atributos podem permitir reidentificação mesmo após a retirada do nome. A divulgação deve considerar esse risco, o consentimento obtido, a finalidade informada e eventuais limites éticos ou jurídicos.

\subsection{Separação entre validação e respostas}

Em coletas digitais, pode ser necessário confirmar elegibilidade ou impedir respostas duplicadas sem associar a identidade ao conteúdo respondido. Uma arquitetura possível é validar o participante em uma etapa separada, emitir um código de uso único e armazenar as respostas sem identificadores diretos. A tabela de correspondência, se indispensável, deve ficar separada, protegida e acessível apenas a responsáveis autorizados.

\section{Validade e autenticidade da coleta}

A proteção da identidade deve conviver com a qualidade da pesquisa. O projeto precisa definir critérios de inclusão, prevenção de duplicidade, teste do instrumento, registro de recusas e perdas e procedimentos para identificar respostas inválidas sem expor participantes.

Se a pesquisa pretende fazer inferências sobre uma população, a estratégia amostral e seus limites devem ser explícitos. Quantidade de respostas não corrige uma amostra enviesada. Controles de autenticidade, quando necessários, devem ser proporcionais ao risco e informados aos participantes; vigilância excessiva pode violar privacidade e comprometer a própria pesquisa.

\section{Ética em pesquisa bibliográfica e documental}

Pesquisas bibliográficas e documentais também exigem ética. O estudante deve representar corretamente as ideias dos autores, evitar distorções, citar as fontes consultadas e distinguir análise própria de informação obtida em documentos.

Fontes digitais devem ser avaliadas com cuidado. Sempre que possível, devem ser priorizados livros, artigos científicos, legislação, documentos institucionais, bases oficiais e publicações técnicas reconhecidas.

\section{Plágio, citação e crédito acadêmico}

Plágio é apresentar texto, ideia, dado, imagem, estrutura ou resultado alheio como se fosse próprio. A cópia literal sem indicação, a paráfrase sem referência e a tradução não identificada omitem a autoria e violam a integridade acadêmica.

Citar não é plagiar. O estudante pode reproduzir um trecho nos limites pertinentes ou explicar uma ideia com suas próprias palavras, desde que identifique a fonte de acordo com as regras de citação e inclua a referência correspondente. A citação dá crédito; ela não autoriza montagem excessiva de texto alheio nem substitui a análise do autor do TCC.

\caixaorientacao{Ao usar um trabalho anterior, deixe visíveis três camadas: o que pertence à fonte, onde a fonte pode ser localizada e qual é a interpretação ou contribuição própria do TCC. O capítulo sobre integridade acadêmica detalha as formas de plágio e sua prevenção.}

Para os tipos de plágio, o uso responsável de tecnologias e as medidas de prevenção, consulte o Capítulo \ref{chap:integridade}.

\section{Conflitos de interesse}

Quando o estudante pesquisa seu próprio local de trabalho, estágio, serviço ou instituição, deve reconhecer possíveis conflitos de interesse e vieses. A proximidade com o objeto pode enriquecer a pesquisa, mas exige transparência, proteção contra coerção e cuidado na interpretação.

\caixaorientacao{A ética precisa aparecer no desenho da pesquisa, na coleta, na análise, na escrita, no armazenamento e na divulgação dos resultados.}

\section{Retenção, descarte e compartilhamento}

O protocolo e o termo de consentimento, quando aplicáveis, devem ser coerentes sobre prazo de guarda, possibilidade de reutilização e compartilhamento. O descarte precisa ser seguro e verificável. Antes de depositar bases junto ao TCC, confirme se a divulgação é autorizada e se os dados estão adequadamente anonimizados.

\section{Três autorizações com funções diferentes}

A autorização da instituição pesquisada permite acesso ao local, aos registros ou à atividade, dentro de seus limites. O consentimento ou assentimento protege a decisão livre do participante. A aprovação ética avalia o protocolo no sistema competente. Um documento não substitui automaticamente os demais.
